% ! TeX root = distributed-systems-final-report.tex
\section{Deployment}\label{deployment}

The system is deployed using docker compose, so a working \textit{Docker Engine}
installation and \textit{docker compose v2} are required. Internet access is
required to pull the two used images from the web: \textit{python:3.12-slim} and
\textit{eclipse-mosquitto:2.1.2-alpine}.

To clone the repository from github:
\begin{verbatim}
git clone \
  https://github.com/papum20/unibo__projects__distributed-smart-roundabout.git
\end{verbatim}

The project is already setup and ready to be immediately deployed; changing
configuration variables is optional. One should just be careful to have the
docker's used ports free on his computer (currently, only 8080, for the web
service).

From linux, it can be \textbf{run} using the available bash scripts, which
should be launched from the project's root directory. To run it:
\begin{verbatim}
./src/start.sh
\end{verbatim}

Details on the running phase can be found in \ref{user-guide}.

To stop docker compose (if necessary, first stop watching the containers logs
by pressing \textit{Ctrl+C}):
\begin{verbatim}
./src/down.sh
\end{verbatim}

Remember that any variable change requires to restart the proper containers; to
be safe, restart all with \textit{down.sh} followed by \textit{start.sh}.

The only \textit{docker compose} file \textit{src/docker-compose.yml} reads
\textbf{environment variables} from
\textit{src/bash-scripts/clean.env.generated}, but they can be configured from
\textit{src/.env} in a more readable way (with spacing): the
\textit{src/start.sh} script is a wrapper that cleans \textit{.env} to produce
\textit{clean.env.generated}, before deploying with docker compose. Variables
are explained through comments in the file and can be better understood by
reading this report, but, in general, they involve the following topics:
\begin{itemize}
  \item ports (for container communication inside the virtual network) and local
  ports (exposed on the machine);
  \item message topics names;
  \item number of vehicle replicas (number of replicated docker containers).
\end{itemize}

Other variables, related to the implementation itself, can be configured in
\textit{common/const.py}: they will be shared among all python containers. The
most interesting ones to try out are:
\begin{itemize}
  \item dimensions for roundabout, roads and visualization area;
  \item number of roads (\textit{ROUNDABOUT\_N\_ROADS});
  \item additional vehicle safety margin \textit{VEHICLE\_SAFETY\_MARGIN\_M} (refer to
  \ref{arch:behaviour_controller});
  \item car's local vision radius \textit{CAR\_VISION\_RADIUS\_M};
\end{itemize}

Other variables concern:
\begin{itemize}
  \item updates per second, for both vehicles and controller;
  \item network timeout;
  \item vehicle dimensions and other technical details;
  \item roundabout proximity distance: threshold to consider a vehicle as
  ``about to enter'';
  \item controller's deadlock timer (as seen in
  \ref{arch:behaviour_controller});
\end{itemize}

All variables try to mimic real measures: in meters for distances, seconds for
time intervals.
\\
\\
\textit{About directories and docker compose structure}: in \textit{src/},
there's a separate folder for each service. Except mosquittoMQTT, which instead
directly uses the standard docker image (with a minimal configuration file), all
other services' images copy their own folder into their respective container and
run python. The \textit{common/} folder contains common variables and utilities
and is copied into all python containers.
